%%%%%%%%%%  scan idx 27  |  folio 14  %%%%%%%%%%
[1] un
homomorphisme groupes
\begin{equation}\tag{5.2}
f^{\mathrm{top}}_*:K^\bullet(X^{\mathrm{top}})\longrightarrow K^\bullet(Y^{\mathrm{top}}).
\end{equation}
Utilisant de m\^eme l'homomorphisme $f_*$ de (3.1) d\'efini par la formule (3.2) (qui utilise ici implicitement le th\'eor\`eme de finitude de Grauert [4] pour un morphisme propre d'espaces analytiques) on trouve un carr\'e d'homomorphismes naturels, ne faisant intervenir que des groupes du type $K^\bullet$ (\`a l'exclusion d'anneaux du type anneaux de Chow ou de cohomologie)~:
\begin{equation}\tag{5.3}
\begin{array}{ccc}
K^\bullet(X) & \longrightarrow & K^\bullet(X^{\mathrm{top}}) \\
f_*\downarrow && \downarrow f^{\mathrm{top}}_* \\
K^\bullet(Y) & \longrightarrow & K^\bullet(Y^{\mathrm{top}}).
\end{array}
\end{equation}
La formule de type Riemann--Roch \`a laquelle nous faisions allusion consiste \`a affirmer la commutativit\'e de ce carr\'e. On notera la diff\'erence de nature entre cette assertion (qui ne n\'eglige pas les ph\'enom\`enes de torsion, et se place dans un groupe $K^\bullet(Y^{\mathrm{top}})$ de nature essentiellement discr\`ete), et la formule de Riemann--Roch du type envisag\'e au n\textsuperscript{o}~4, se pla\c{c}ant dans un groupe du type «~Chow~» qui n'est plus de nature discr\`ete, mais qu'on doit tensoriser en revanche par $\mathbb Q$. Notons que la formule (5.3) est d\'emontr\'ee en tous cas lorsque $X$ et $Y$ sont des vari\'et\'es projectives non singuli\`eres [7], comme la formule du type envisag\'e au n\textsuperscript{o}~4, qui est d\'emontr\'ee dans ce cas dans [RRR]. Il s'agit l\`a d'une variante plausible du th\'eor\`eme de l'index de Atiyah--Singer (et de ses g\'en\'eralisations dues \`a Shih [7] et Atiyah), qui appelle \'evidemment une g\'en\'eralisation commune. Une fois prouv\'ee la commutativit\'e de (5.3), on en conclurait aussit\^ot, \`a l'aide du «~th\'eor\`eme de Riemann--Roch diff\'erentiable~» de Atiyah--Hirzebruch [1], une variante cohomologique de la formule de Riemann--Roch analytique complexe, savoir la commutativit\'e du carr\'e
\begin{equation}\tag{5.4}
\begin{array}{ccc}
K^\bullet(X) & \xrightarrow{\;x\longrightarrow \ch(x)\Todd(T_X)\;} & H^{2*}(X,\mathbb Q) \\
\downarrow && \downarrow \\
K^\bullet(Y) & \xrightarrow{\;y\longrightarrow \ch(y)\Todd(T_Y)\;} & H^{2*}(Y,\mathbb Q),
\end{array}
\end{equation}
o\`u les fl\^eches horizontales sont d\'eduites des classes de Chern cohomologiques, d\'eduites de l'homomorphisme (5.1) et de l'homomorphisme classe de Chern ordinaire
\[
K^\bullet(X^{\mathrm{top}})\longrightarrow H^{2*}(X,\mathbb Z),
\]
